\documentclass[10pt,a4paper]{amsart}
\usepackage[margin=19mm]{geometry}
\usepackage{amsmath,amssymb,mathtools}
\usepackage[scr=rsfs,cal=euler]{mathalfa}
\usepackage{xfrac}
\usepackage[unicode,hidelinks,bookmarks=false]{hyperref}
\newcommand{\ie}{i.e.\ }
\newcommand{\cf}{cf.\ }
\newcommand{\resp}{respectively\ }
\newcommand{\Subsection}{\subsection}
\providecommand{\nobreakdash}{\nobreak\hbox{-}}
\setlength{\emergencystretch}{2em}
\multlinegap0pt
\usepackage[shortlabels]{enumitem}
\usepackage{amssymb}
\newtheorem{lemma}{Lemma}
\newcommand{\eps}{\varepsilon}
\newcommand{\pr}{\mathrm{Pr}}
\title{A Question of Erd\H{o}s and Graham on Egyptian Fractions}
\author{David Conlon, Jacob Fox, Xiaoyu He, Dhruv Mubayi, Huy Tuan Pham, Andrew Suk, Jacques Verstra\"ete}
\date{Source: arXiv:2404.16016v2 (CC BY 4.0)}
\begin{document}
\maketitle

\noindent\emph{Source and licence.} A Question of Erd\H{o}s and Graham on Egyptian Fractions. Version arXiv:2404.16016v2 (CC BY 4.0), distributed by arXiv under the Creative Commons Attribution 4.0 International licence, http://creativecommons.org/licenses/by/4.0/, as recorded in the arXiv licence metadata for this exact version and shown on the abstract page https://arxiv.org/abs/2404.16016v2. Full licence notice: \texttt{licenses/arxiv-2404.16016-CC-BY-4.0.txt}.\par\smallskip

\noindent\textit{Editorial scope.} Counted unit: the article's lemma lem:main-ent (source0.tex lines 155--158). The article's complete original proof of it is delivered with it (source0.tex lines 219--309, one \texttt{proof} environment of 91 lines, which the article places away from the statement). No mathematics is written, repaired or re-derived: the statement and the proof are the article's own lines, in the article's own notation, and no retained line is substituted or re-flowed.  Retained uncounted as the source-local context the proof needs: lines 162--174; lines 212--217.  Nothing was omitted from any retained span, so the package carries no omission record.  No retained line contains a citation, so the wrapper carries no bibliography and no bibliography entry.  No retained line contains a figure environment, an image-inclusion command or drawing code, so no image is delivered and the package declares no asset dependency.  Editorial formatting only, in addition: an AMS \texttt{amsart} preamble that restates the article's own declarations and macros used by the retained lines, namely the environments \texttt{lemma} on separate counters, together with the standard packages those lines need.  The retained environments print in this document as Lemma 1 in order of appearance, whereas the article prints the same items with the numbering of its own full document.  The complete native source of the article as distributed in the arXiv e-print for this exact version is bundled unchanged as \texttt{sources/158107/source0.tex}.
\begin{lemma}\label{lem:opt-dis}
    For $x < (1/2)(\sum_{m=1}^{n} 1/m)$, the distribution $P(x)$ is given by setting $p_m = \frac{1}{1+e^{c_{x,n} n/m}}$ for the unique $c_{x,n}$ such that $\sum_{m=1}^{n} p_m/m = x$. 
    Moreover, there is $C>0$ such that, provided $x < (1-\eps)(\ln n)/2$, $c_{x,n} > 0$ and $c_{x,n} \in [C^{-1}e^{-2x}, Ce^{-2x}]$. 
    Finally, for $\eps>0$ and assuming $x\in [\eps,1/\eps]$, $c_{x,n} = \lambda+o_n(1)$, where $\lambda > 0$ is the unique solution to 
    \[
        \int_{0}^{1} \frac{1}{y(1+e^{\lambda/y})} dy = x.
    \]
    In particular, 
    \begin{equation}\label{eq:ent-cal}
        H(P(x)) = (1+o_n(1)) n \int_{0}^{1} h\left(\frac{1}{1+e^{\lambda/y}}\right)dy = \Theta(n).
    \end{equation}
    Note that here the $o_n(1)$ terms may depend on $\eps$. 
\end{lemma}

\begin{lemma}\label{lem:B-E}
    Let $X_1,\dots,X_n$ be independent centered random variables with $\mathbb{E}[X_i^2] = \zeta_i$ and $\mathbb{E}[|X_i|^3] = \rho_i$. Let $Z = \frac{\sum_{i=1}^{n}X_i}{(\sum_{i=1}^{n}\zeta_i)^{1/2}}$ and $Z'$ be a standard Gaussian. Then 
    \[
        \sup_{y\in \mathbb{R}} |\pr[Z \le y] - \pr[Z' \le y]| \le C\frac{\sum_{i=1}^{n}\rho_i}{(\sum_{i=1}^{n}\zeta_i)^{3/2}}.
    \]
\end{lemma}

\begin{lemma}\label{lem:main-ent}
    For $\eps>0$ and $x \in (0,(1-\eps)(\ln n)/2)$, the number of subsets $A\subseteq [n]$ with $s(A) \le x$ is bounded above by $2^{H(P(x))}$. 
    Furthermore, there is $c_{x,n}>0$ with $c_{x,n} = \Theta(e^{-2x})$ such that, for any $U \subseteq [n]$, the number of subsets $A\subseteq U$ with $s(A) \le x$ is bounded below by $2^{H(P(x)) - (n-|U|) - O(\sqrt{n/c_{x,n}})}$.
\end{lemma}

\begin{proof}[Proof of Lemma \ref{lem:main-ent}]
Let $S$ be a finite set of real numbers and $n=|S|$. Let $r_S(x)$ be the number of subsets of $S$ that sum to at most $x$. Define the random variable $X$ to be a uniform random subset of $S$ whose elements sum to at most $x$. Note that the entropy of $X$ satisfies $H(X)=\log_2 r_S(x)$, so $r_S(x)=2^{H(X)}$. For each $s \in S$, let $X_s$ be the indicator random variable of the event $s \in X$ and let $p_s = \Pr[X_s]$, so that $\mathbb{E}[\sum_{s\in S}s X_s]=\sum_{s \in S} sp_s \leq x$. Observe that $X$ has the same distribution as the joint distribution of the $n$ random variables $X_s$. Therefore, by subadditivity of the entropy function, we have 
\[
    H(X) \leq \sum_{s \in S}H(X_s)=\sum_{s \in S}h(p_s).
\]
Hence, we get the upper bound 
$$r_S(x) \leq 2^h,$$ 
where $h$ is the maximum value of $\sum_{s \in S} h(p_s)$ over all choices of $(p_s)_{s \in S}$ satisfying $\sum_{s \in S} sp_s \leq x$. In particular, for $S=\{1/m:m\in [n]\}$, we obtain that the number of subsets $A\subseteq [n]$ with $s(A)\le x$ is at most $2^{H(P(x))}$, as claimed. 


We now turn to the lower bound. Consider independent Bernoulli random variables $Y_m$ for $m \in [n]$ satisfying $Y_m=1$ with probability $p_m$ and $Y_m=0$ otherwise. Let $Y=(Y_m)_{m\in [n]}$, $Z=\sum_{m\in [n]} Y_m/m$ and $E$ be the indicator of the event $Z \le x$. Recall that, for $a\in \{0,1\}$, the conditional entropy $H(Y | E=a) = -\sum_{y\in \{0,1\}^n} \pr[Y=y | E=a] \log \pr[Y=y | E=a]$ and $H(Y | E) = \sum_{a\in \{0,1\}} \pr[E=a] H(Y | E=a)$. Since $E$ is determined by $Y$, we have 
\[
    H(Y) = H(Y,E) = H(Y|E) + H(E) = H(E) + \sum_{a\in \{0,1\}} \pr[E=a] H(Y | E=a).
\]
We thus have
\begin{equation}\label{eq:ent-decomp}
H(Y | E=1) = \frac{1}{\pr[Z\le x]}\left(H(Y) - H(E) - \Pr[Z>x] H(Y | E=0)\right). 
\end{equation}
Let $c=c_{x,n}$ as in Lemma \ref{lem:opt-dis}. By that lemma, the random variable $Z$ has variance 
\begin{align}\label{eq:2nd-mm}
    \sum_{m=1}^{n} \left(\frac{1}{(1+e^{cn/m})m^2} - \frac{1}{(1+e^{cn/m})^2m^2}\right) = \Theta\left(\sum_{m=1}^{n} \frac{e^{-cn/m}}{m^2}\right) = \Theta\left(\frac{1}{cn}\right). 
\end{align}
To see the last bound, observe that, for $k \ge 1$,  
$\sum_{m=cn/(k+1)}^{cn/k} \frac{e^{-cn/m}}{m^2} = \Theta(\frac{e^{-k}}{cn})$. Similarly, 
for $1\le k\le 1/c$, $\sum_{m=cnk}^{cn(k+1)} \frac{e^{-cn/m}}{m^2} = \Theta(\frac{e^{-1/k}}{cnk^2})$. The desired bound follows from summing these estimates over $k$. 

By a similar argument, the sum of the absolute centered third moments of the $Y_m/m$ is 
\begin{align}\label{eq:3rd-mm}
    \sum_{m=1}^{n} \left[\frac{1}{1+e^{cn/m}} \left|\frac{1}{m}-\frac{1}{m(1+e^{cn/m})}\right|^3 + \left(1-\frac{1}{1+e^{cn/m}}\right) \left|-\frac{1}{m(1+e^{cn/m})}\right|^3 \right] &= O\left(\sum_{m=1}^{n} \frac{e^{-cn/m}}{m^3}\right) \nonumber\\
    &= O\left(\frac{1}{(cn)^2}\right). 
\end{align}
The Berry--Esseen bound, Lemma \ref{lem:B-E}, then yields that, for $g\sim \mathcal{N}(0,1)$, 
\begin{equation}\label{eq:pZ}
    \pr(E=1) = \pr[Z\le x] = \pr[g\le 0] + O((cn)^{-1/2}) = 1/2 + O((cn)^{-1/2}),
\end{equation}
where we used that $\mathbb{E}[Z]=x$, together with (\ref{eq:2nd-mm}) and (\ref{eq:3rd-mm}).

We next bound $H(Y_1,\dots,Y_n | E=0) \le \sum_{m=1}^{n} H(Y_m|E=0)$. To bound the summands, we note by Bayes' rule that  \[\pr(Y_m = 1|E=0) = \frac{\pr(E=0|Y_m=1)\pr(Y_m=1)}{\pr(E=0)}\]
and we will use a similar argument with the Berry--Esseen bound to show that $\pr(Y_m = 1|E=0)$ is close to $\pr(Y_m=1)$. 
Indeed, the calculations above similarly yield that the random variable $Z'_m=Z-Y_m/m+1/m$ is a sum of independent random variables with $\mathbb{E}Z'_m = x + \frac{1}{m} - \frac{1}{m(1+e^{cn/m})}$, $\mathrm{Var}(Z'_m) = \Theta(1/(cn))$ and the sum of absolute centered third moments $O(1/(cn)^2)$. 
By Lemma \ref{lem:B-E}, for $g\sim \mathcal{N}(0,1)$, 
\begin{align*}
    \pr(E = 0 | Y_m=1) &= \pr(Z'_m > x) \\
    &= \pr\left(g > -\frac{1/m-1/(m(1+e^{cn/m}))}{\mathrm{Var}(Z'_m)^{1/2}}\right) + O((cn)^{-1/2}) \\
    &= 1/2 + O\left((cn)^{-1/2}+\frac{\sqrt{cn}}{m}\right),
\end{align*}
assuming that $m > 10\sqrt{cn}$ for the last bound, where we used the simple estimate $\pr(g > z) = \frac{1}{2} + O(z)$ for $|z| \le 1$. 
Therefore,
\begin{equation*}
    \left|\frac{\pr(E=0|Y_m=1)}{\pr(E=0)} - 1 \right| = \left|\frac{1/2 + O\left((cn)^{-1/2}+\frac{\sqrt{cn}}{m}\right)}{1/2 + O((cn)^{-1/2})} - 1 \right| \le O\left(\frac{\sqrt{cn}}{m} + \frac{1}{\sqrt{cn}}\right). 
\end{equation*}
Thus, by Bayes' rule,
\begin{equation}\label{eq:prob-bound-1}
    \frac{|\pr(Y_m=1|E=0) - \pr(Y_m=1)|}{\pr(Y_m=1)} \le O\left(\frac{\sqrt{cn}}{m} + \frac{1}{\sqrt{cn}}\right). 
\end{equation}

From (\ref{eq:prob-bound-1}), we have
\[
    \pr(Y_m=1) \left(1 - O\left(\frac{\sqrt{cn}}{m} + \frac{1}{\sqrt{cn}}\right)\right)\le \pr(Y_m=1|E=0) \le \pr(Y_m=1) \left(1 + O\left(\frac{\sqrt{cn}}{m} + \frac{1}{\sqrt{cn}}\right)\right).
\]
Since $h'(p) = \log \frac{1-p}{p} \le \log \frac 1p$, we have that 
\begin{align*}
    &h(\pr(Y_m=1|E=0)) \\
    &\le h(\pr(Y_m=1)) + \left (\log \frac{1}{\pr(Y_m=1) \left(1 - O\left(\frac{\sqrt{cn}}{m} + \frac{1}{\sqrt{cn}}\right)\right)}\right ) \left | \pr(Y_m=1|E=0)-\pr(Y_m=1)\right |\\
    &\le h(\pr(Y_m=1)) + \left (\log \frac{1}{\pr(Y_m=1) \left(1 - O\left(\frac{\sqrt{cn}}{m} + \frac{1}{\sqrt{cn}}\right)\right)}\right ) \cdot O\left(\frac{\sqrt{cn}}{m} + \frac{1}{\sqrt{cn}}\right) \pr(Y_m=1) \\
    &\le h(\pr(Y_m=1)) +  O\left(\frac{\sqrt{cn}}{m} + \frac{1}{\sqrt{cn}}\right) \pr(Y_m=1) \log \frac{1}{\pr(Y_m=1)},
\end{align*}
where in the last inequality we used $\pr(Y_m=1) \le 1/2$, so $\log \frac{1}{\pr(Y_m=1) \left(1 - O\left(\frac{\sqrt{cn}}{m} + \frac{1}{\sqrt{cn}}\right)\right)} = O\left(\log \frac{1}{\pr(Y_m=1)}\right)$. 
Therefore, 
\begin{align*}
    H(Y_1,\dots,Y_n|E=0) &\le \sum_{m=1}^{n} H(Y_m | E=0)\\
    &\le 10\sqrt{cn} + \sum_{m>10\sqrt{cn}} h(\pr(Y_m=1|E=0)) \\
    &\le 10\sqrt{cn} + \sum_{m>10\sqrt{cn}} \left(H(Y_m) + O\left(\frac{\sqrt{cn}}{m} + \frac{1}{\sqrt{cn}}\right) \pr(Y_m=1) \log \frac{1}{\pr(Y_m=1)}\right)\\
    &\le H(Y_1,\dots,Y_n) + 10\sqrt{cn} + \sum_{m>10\sqrt{cn}}O\left(\frac{\sqrt{cn}}{m} + \frac{1}{\sqrt{cn}}\right) \pr(Y_m=1) \log \frac{1}{\pr(Y_m=1)}\\
    &\le H(Y_1,\dots,Y_n) + 10\sqrt{cn} + \sum_{m>10\sqrt{cn}} O\left(\frac{\sqrt{cn}}{m} + \frac{1}{\sqrt{cn}}\right) \frac{cn/m}{1+e^{cn/m}}\\
    &\le H(Y_1,\dots,Y_n) + O(\sqrt{n/c}). 
\end{align*}
Again, for the last estimate, we note that, for $k \ge 1$, $\sum_{m=cn/(k+1)}^{cn/k} \left(\frac{\sqrt{cn}}{m} + \frac{1}{\sqrt{cn}}\right) \frac{cn/m}{1+e^{cn/m}} \le O\left(\frac{cn}{k^2} e^{-k}\frac{k^2}{\sqrt{cn}}\right) = O(e^{-k}\sqrt{cn})$  and, for $1\le k\le 1/c$, $\sum_{m=cnk}^{cn(k+1)} \left(\frac{\sqrt{cn}}{m} + \frac{1}{\sqrt{cn}}\right) \frac{cn/m}{1+e^{cn/m}} \le O\left(cn\frac{1}{k\sqrt{cn}}\right) \le O\left(\frac{1}{k}\sqrt{cn}\right)$. Summing over $k$, we thus have 
\[
    \sum_{m>10\sqrt{cn}} O\left(\frac{\sqrt{cn}}{m} + \frac{1}{\sqrt{cn}}\right) \frac{cn/m}{1+e^{cn/m}} = O(\sqrt{n/c}).
\]
Combining with (\ref{eq:ent-decomp}), and noting that $H(E)\le 1$ and $\pr(Z\le x) = 1/2 + O((cn)^{-1/2})$ by (\ref{eq:pZ}), we obtain that 
\[
    H(Y_1,\dots,Y_n|E=1) \ge (1-O((cn)^{-1/2}))(H(P(x)) - O(\sqrt{n/c})) = H(P(x)) - O(\sqrt{n/c}).
\]
Using that for any random variable $X$ we have $H(X)\le \log |\mathrm{supp}(X)|$, we obtain that the number of subsets $A\subseteq [n]$ with $s(A) \le x$ is at least $2^{H(P(x)) - O(\sqrt{n/c})}$. This implies that the number of subsets $A\subseteq [n]\setminus U$ with $s(A)\le x$ is at least 
\[
    2^{-(n-|U|)} 2^{H(P(x)) - O(\sqrt{n/c})},
\]
as required. 
\end{proof}

\end{document}
